% ============================================================================
% DM-2026-008: DI-PSI-CDH-GDS Desk Instruction CUI Classification
% ============================================================================
\newcommand{\UniqueID}{DM-2026-008}
\newcommand{\DocumentDate}{April 13, 2026}
\newcommand{\AuthorName}{Bruce Dombrowski}
\newcommand{\AuthorTitle}{Systems Engineer, Payload Software Integration}
\newcommand{\ToField}{PSI Team, PIM, NASA JSC CUI Program Manager}
\newcommand{\SubjectField}{CUI Classification: DI-PSI-CDH-GDS Desk Instruction}
\newcommand{\OPRField}{Payload Software Integration (PSI)}

\newcommand{\dmContent}{%
\begin{enumerate}

\dmsection{Purpose}

This memorandum documents the decision to classify the DI-PSI-CDH-GDS ``C\&DH / GDS / IT Security Coordination Meeting Guide'' desk instruction as \textbf{CUI//SP-EXPT//FEDCON} in its as-shipped state, and to flag the potential applicability of the \textbf{ISVI (Information Systems Vulnerability Information)} category for NASA JSC CUI Program Manager adjudication.

\dmsection{Background}

DI-PSI-CDH-GDS is a 25-page companion desk instruction for the Doc 8604 Rev D checklist. Unlike the blank checklist (which contains only empty form fields and publicly-citable SSP document numbers), the desk instruction contains \textbf{procedural knowledge about ISS information infrastructure}:

\begin{itemize}
    \item \textbf{NTP architecture:} iPEHG-2 ISS/GPS time path and SSC NTP server ground-time path
    \item \textbf{Antivirus update distribution:} PL NAS, SSC System, and PD-uplink methods with access protocol specifics (SMB, NFS/AFP, HTTP URL)
    \item \textbf{KuIP protocol assignment rules:} Port-numbering constraints ($>$1023), APID encoding structure, NAT IP assignment procedures, Space/Ground Node ID management
    \item \textbf{JSL network topology:} Hirschmann Switch, WAPs, Edge Router connection types, static route handling, VLAN configuration guidance
    \item \textbf{HOSC GDS routing:} PDSS CCSDS Space Packet routing, iVODS paths, ICU on-board LOS storage, GSE packet assignment
    \item \textbf{D684-11372-01 identifier encoding:} APID bit-field structure, Subset ID encoding rules, PUI field definitions, KuIP PDL Log management procedures
    \item \textbf{MAC address assignment procedures:} Configurable vs.\ hardcoded determination, who assigns (PSI vs.\ PD), and how assignments flow to the Preliminary SW ICD
    \item \textbf{IT Security:} SSP 50974 requirements flow, malware-definition update paths, SSH key pair and X.509 certificate exchange procedures
\end{itemize}

This procedural knowledge describes \textit{how the ISS information infrastructure operates}, even without disclosing specific IP addresses or payload configurations. An adversary with this knowledge would understand the architecture, data paths, assignment authority chains, and identifier structures of the ISS C\&DH system.

\dmsection{Analysis}

\textbf{Export Controlled (SP-EXPT):} The desk instruction describes technical data about the ISS payload C\&DH and GDS systems that is subject to U.S. export control under the EAR (15 CFR Parts 730--774) and/or ITAR (22 CFR Parts 120--130). The ISS C\&DH system architecture, including protocol assignment rules, identifier encoding, and network topology, constitutes controlled technical data about a space system. \textbf{SP-EXPT applies.}

\textbf{Information Systems Vulnerability Information (ISVI):} The CUI Registry defines ISVI as information about vulnerabilities in information systems. The desk instruction describes:
\begin{itemize}
    \item Network topology (which devices connect where)
    \item Authentication paths (SSH key exchange, X.509 certificate procedures)
    \item Time-synchronization architecture (NTP sources)
    \item Malware-definition update distribution channels
    \item Identifier encoding that could aid spoofing (APID bit-field structure)
\end{itemize}

Whether this constitutes ``vulnerability information'' or merely ``system description'' is a judgment call that \textbf{requires NASA JSC CUI PM adjudication}. This decision conservatively marks the document as SP-EXPT and flags ISVI for determination.

\textbf{Limited Dissemination Control --- FEDCON:} The desk instruction is used by Boeing PSI engineers (federal contractors) supporting NASA payloads. The ISS is an international-partner program (JAXA, ESA, CSA, Roscosmos), so NOFORN is inappropriate. FEDCON (Federal Employees and Contractors Only) is the correct LDC. If international-partner access is needed, the NASA JSC CUI PM may authorize REL TO markings on a case-by-case basis.

\dmsection{Decision}

\textbf{DI-PSI-CDH-GDS is classified as CUI//SP-EXPT//FEDCON.}

\begin{itemize}
    \item The desk instruction carries CUI//SP-EXPT//FEDCON banner on every page (running header and footer)
    \item The cover page explicitly states the classification and notes the ISVI-pending determination
    \item The desk instruction must be distributed only via NASA ORBIT Files (preferred) or encrypted email (DoD SAFE, NASA Box CUI, FIPS 140-2/3 S/MIME)
    \item The desk instruction must not be committed to any public repository
    \item The desk instruction must not be emailed unencrypted
    \item Destroy copies per NIST SP 800-88
\end{itemize}

\textbf{Pending action:} NASA JSC CUI PM to adjudicate whether ISVI additionally applies. If yes, the banner becomes \textbf{CUI//SP-EXPT/ISVI//FEDCON}. This is tracked as compliance finding F-12 in \texttt{docs/compliance/findings-and-recommendations.md}.

\textbf{Distinction from the blank checklist:} Per DM-2026-007, the blank Doc 8604 Rev D checklist form is NOT CUI because it contains only empty fields and publicly-citable document numbers. The desk instruction is CUI because it contains the \textit{procedural answers} --- the infrastructure knowledge that the blank form merely asks about.

\dmsection{Implementation}

\begin{itemize}
    \item CUI//SP-EXPT//FEDCON banner added to DI-PSI-CDH-GDS-Meeting-Guide.pdf running header/footer via \texttt{docs/di-pandoc-header.tex}
    \item Cover page classification block updated via \texttt{docs/di-pandoc-before-body.tex}
    \item ISVI-pending note included on cover page referencing compliance finding F-12
    \item Repository (\texttt{ISS\_Payload\_CDH\_Interface\_PSI\_Checklist}) is PRIVATE on GitHub --- verified, no public-release violation
    \item The companion PSI Training Module spreadsheets (\texttt{PSI-TrainingModule\_600-601*.xltx}, \texttt{PSI-TrainingModule\_602*.xlsx}) are also CUI//SP-EXPT and co-reside in the same private repository
\end{itemize}

\dmsection{References}

\begin{itemize}
    \item 32 CFR Part 2002.4 (CUI definition): \url{https://www.law.cornell.edu/cfr/text/32/part-2002}
    \item 32 CFR Part 2002.20 (CUI marking requirements): \url{https://www.law.cornell.edu/cfr/text/32/section-2002.20}
    \item NARA CUI Registry --- Export Controlled: \url{https://www.archives.gov/cui/registry/category-detail/export-control}
    \item NARA CUI Registry --- ISVI: \url{https://www.archives.gov/cui/registry/category-detail/information-systems-vulnerability-information}
    \item NPR 2810.7 ``Controlled Unclassified Information'': \url{https://nodis3.gsfc.nasa.gov/displayDir.cfm?t=NPR&c=2810&s=7}
    \item \texttt{docs/compliance/findings-and-recommendations.md} --- Finding F-12 (ISVI dual-category question)
    \item \texttt{docs/compliance/blank-form-cui-determination.md} --- blank-form NOT CUI determination
    \item \texttt{docs/compliance/correctness-review.md} --- citation verification
    \item DM-2026-006 (NASA CUI Directive Requirements Absorption)
    \item DM-2026-007 (Blank Checklist NOT CUI Determination)
\end{itemize}

\end{enumerate}
}

% ============================================================================
% Decision Memorandum Base Template
% ============================================================================
% This is the shared base template for all Decision Memoranda.
% DO NOT edit this file for individual decisions.
%
% Usage: Create a thin wrapper that defines variables and content, then:
%   % ============================================================================
% Decision Memorandum Base Template
% ============================================================================
% This is the shared base template for all Decision Memoranda.
% DO NOT edit this file for individual decisions.
%
% Usage: Create a thin wrapper that defines variables and content, then:
%   % ============================================================================
% Decision Memorandum Base Template
% ============================================================================
% This is the shared base template for all Decision Memoranda.
% DO NOT edit this file for individual decisions.
%
% Usage: Create a thin wrapper that defines variables and content, then:
%   \input{_template.tex}
%
% Required variables (define before \input):
%   \UniqueID       - e.g., DM-2026-001
%   \DocumentDate   - e.g., January 19, 2026
%   \AuthorName     - Signer name
%   \AuthorTitle    - Signer title
%   \ToField        - Distribution
%   \SubjectField   - Subject line
%   \OPRField       - Office of Primary Responsibility
%   \dmContent      - The full content (enumerate environment with \dmsection items)
%
% ============================================================================

\documentclass[11pt,letterpaper]{article}

% ============================================================================
% PACKAGES
% ============================================================================
\usepackage[margin=1in, top=1.15in, headheight=30pt]{geometry}
\usepackage{graphicx}
\usepackage{fancyhdr}
\usepackage{lastpage}
\usepackage{enumitem}
\usepackage{xcolor}
\usepackage{hyperref}
\usepackage{tabularx}

% ============================================================================
% DOCUMENT CONFIGURATION
% ============================================================================
\hypersetup{
    colorlinks=true,
    linkcolor=blue,
    urlcolor=blue,
    pdfborder={0 0 0}
}

% Remove paragraph indentation
\setlength{\parindent}{0pt}
\setlength{\parskip}{6pt}

% List formatting - match Word style
\setlist[enumerate]{leftmargin=0.5in, labelwidth=0.25in, labelsep=0.1in, align=left, itemsep=12pt, topsep=12pt}
\setlist[enumerate,1]{label=\arabic*.}
\setlist[itemize]{leftmargin=0.5in, itemsep=3pt, topsep=6pt}

% ============================================================================
% HEADER AND FOOTER CONFIGURATION
% ============================================================================
\pagestyle{fancy}
\fancyhf{}

% Header: text-only, Boeing / ISS Program identification
\fancyhead[L]{\large\textbf{Decision Memorandum}\\[-2pt]\scriptsize Boeing \textbar{} International Space Station Program}
\fancyhead[R]{\small \UniqueID\\[-2pt]\scriptsize \DocumentDate}
\renewcommand{\headrulewidth}{0.4pt}

% Footer: Unique ID (left), Page X of Y (center), Date (right)
\fancyfoot[L]{\small \UniqueID}
\fancyfoot[C]{\small Page \thepage\ of \pageref{LastPage}}
\fancyfoot[R]{\small \DocumentDate}
\renewcommand{\footrulewidth}{0.4pt}

% ============================================================================
% CUSTOM COMMANDS
% ============================================================================
% Bold section headers for numbered sections
\newcommand{\dmsection}[1]{\item \textbf{#1}\par}

% ============================================================================
% BEGIN DOCUMENT
% ============================================================================
\begin{document}

% ----------------------------------------------------------------------------
% UNIQUE ID (top right)
% ----------------------------------------------------------------------------
\hfill \UniqueID

\vspace{0.25in}

% ----------------------------------------------------------------------------
% SIGNATURE BLOCK
% ----------------------------------------------------------------------------
\underline{\hspace{2.5in}}\\[2pt]
\AuthorName\\
\AuthorTitle

\vspace{0.25in}

% ----------------------------------------------------------------------------
% MEMO HEADER FIELDS
% ----------------------------------------------------------------------------
\textbf{DATE:} \DocumentDate

\textbf{TO:} \ToField

\textbf{SUBJECT:} \SubjectField

\textbf{OFFICE OF PRIMARY RESPONSIBILITY:} \OPRField

\vspace{0.25in}

% ============================================================================
% MAIN CONTENT - NUMBERED SECTIONS
% ============================================================================
\dmContent

\end{document}

%
% Required variables (define before \input):
%   \UniqueID       - e.g., DM-2026-001
%   \DocumentDate   - e.g., January 19, 2026
%   \AuthorName     - Signer name
%   \AuthorTitle    - Signer title
%   \ToField        - Distribution
%   \SubjectField   - Subject line
%   \OPRField       - Office of Primary Responsibility
%   \dmContent      - The full content (enumerate environment with \dmsection items)
%
% ============================================================================

\documentclass[11pt,letterpaper]{article}

% ============================================================================
% PACKAGES
% ============================================================================
\usepackage[margin=1in, top=1.15in, headheight=30pt]{geometry}
\usepackage{graphicx}
\usepackage{fancyhdr}
\usepackage{lastpage}
\usepackage{enumitem}
\usepackage{xcolor}
\usepackage{hyperref}
\usepackage{tabularx}

% ============================================================================
% DOCUMENT CONFIGURATION
% ============================================================================
\hypersetup{
    colorlinks=true,
    linkcolor=blue,
    urlcolor=blue,
    pdfborder={0 0 0}
}

% Remove paragraph indentation
\setlength{\parindent}{0pt}
\setlength{\parskip}{6pt}

% List formatting - match Word style
\setlist[enumerate]{leftmargin=0.5in, labelwidth=0.25in, labelsep=0.1in, align=left, itemsep=12pt, topsep=12pt}
\setlist[enumerate,1]{label=\arabic*.}
\setlist[itemize]{leftmargin=0.5in, itemsep=3pt, topsep=6pt}

% ============================================================================
% HEADER AND FOOTER CONFIGURATION
% ============================================================================
\pagestyle{fancy}
\fancyhf{}

% Header: text-only, Boeing / ISS Program identification
\fancyhead[L]{\large\textbf{Decision Memorandum}\\[-2pt]\scriptsize Boeing \textbar{} International Space Station Program}
\fancyhead[R]{\small \UniqueID\\[-2pt]\scriptsize \DocumentDate}
\renewcommand{\headrulewidth}{0.4pt}

% Footer: Unique ID (left), Page X of Y (center), Date (right)
\fancyfoot[L]{\small \UniqueID}
\fancyfoot[C]{\small Page \thepage\ of \pageref{LastPage}}
\fancyfoot[R]{\small \DocumentDate}
\renewcommand{\footrulewidth}{0.4pt}

% ============================================================================
% CUSTOM COMMANDS
% ============================================================================
% Bold section headers for numbered sections
\newcommand{\dmsection}[1]{\item \textbf{#1}\par}

% ============================================================================
% BEGIN DOCUMENT
% ============================================================================
\begin{document}

% ----------------------------------------------------------------------------
% UNIQUE ID (top right)
% ----------------------------------------------------------------------------
\hfill \UniqueID

\vspace{0.25in}

% ----------------------------------------------------------------------------
% SIGNATURE BLOCK
% ----------------------------------------------------------------------------
\underline{\hspace{2.5in}}\\[2pt]
\AuthorName\\
\AuthorTitle

\vspace{0.25in}

% ----------------------------------------------------------------------------
% MEMO HEADER FIELDS
% ----------------------------------------------------------------------------
\textbf{DATE:} \DocumentDate

\textbf{TO:} \ToField

\textbf{SUBJECT:} \SubjectField

\textbf{OFFICE OF PRIMARY RESPONSIBILITY:} \OPRField

\vspace{0.25in}

% ============================================================================
% MAIN CONTENT - NUMBERED SECTIONS
% ============================================================================
\dmContent

\end{document}

%
% Required variables (define before \input):
%   \UniqueID       - e.g., DM-2026-001
%   \DocumentDate   - e.g., January 19, 2026
%   \AuthorName     - Signer name
%   \AuthorTitle    - Signer title
%   \ToField        - Distribution
%   \SubjectField   - Subject line
%   \OPRField       - Office of Primary Responsibility
%   \dmContent      - The full content (enumerate environment with \dmsection items)
%
% ============================================================================

\documentclass[11pt,letterpaper]{article}

% ============================================================================
% PACKAGES
% ============================================================================
\usepackage[margin=1in, top=1.15in, headheight=30pt]{geometry}
\usepackage{graphicx}
\usepackage{fancyhdr}
\usepackage{lastpage}
\usepackage{enumitem}
\usepackage{xcolor}
\usepackage{hyperref}
\usepackage{tabularx}

% ============================================================================
% DOCUMENT CONFIGURATION
% ============================================================================
\hypersetup{
    colorlinks=true,
    linkcolor=blue,
    urlcolor=blue,
    pdfborder={0 0 0}
}

% Remove paragraph indentation
\setlength{\parindent}{0pt}
\setlength{\parskip}{6pt}

% List formatting - match Word style
\setlist[enumerate]{leftmargin=0.5in, labelwidth=0.25in, labelsep=0.1in, align=left, itemsep=12pt, topsep=12pt}
\setlist[enumerate,1]{label=\arabic*.}
\setlist[itemize]{leftmargin=0.5in, itemsep=3pt, topsep=6pt}

% ============================================================================
% HEADER AND FOOTER CONFIGURATION
% ============================================================================
\pagestyle{fancy}
\fancyhf{}

% Header: text-only, Boeing / ISS Program identification
\fancyhead[L]{\large\textbf{Decision Memorandum}\\[-2pt]\scriptsize Boeing \textbar{} International Space Station Program}
\fancyhead[R]{\small \UniqueID\\[-2pt]\scriptsize \DocumentDate}
\renewcommand{\headrulewidth}{0.4pt}

% Footer: Unique ID (left), Page X of Y (center), Date (right)
\fancyfoot[L]{\small \UniqueID}
\fancyfoot[C]{\small Page \thepage\ of \pageref{LastPage}}
\fancyfoot[R]{\small \DocumentDate}
\renewcommand{\footrulewidth}{0.4pt}

% ============================================================================
% CUSTOM COMMANDS
% ============================================================================
% Bold section headers for numbered sections
\newcommand{\dmsection}[1]{\item \textbf{#1}\par}

% ============================================================================
% BEGIN DOCUMENT
% ============================================================================
\begin{document}

% ----------------------------------------------------------------------------
% UNIQUE ID (top right)
% ----------------------------------------------------------------------------
\hfill \UniqueID

\vspace{0.25in}

% ----------------------------------------------------------------------------
% SIGNATURE BLOCK
% ----------------------------------------------------------------------------
\underline{\hspace{2.5in}}\\[2pt]
\AuthorName\\
\AuthorTitle

\vspace{0.25in}

% ----------------------------------------------------------------------------
% MEMO HEADER FIELDS
% ----------------------------------------------------------------------------
\textbf{DATE:} \DocumentDate

\textbf{TO:} \ToField

\textbf{SUBJECT:} \SubjectField

\textbf{OFFICE OF PRIMARY RESPONSIBILITY:} \OPRField

\vspace{0.25in}

% ============================================================================
% MAIN CONTENT - NUMBERED SECTIONS
% ============================================================================
\dmContent

\end{document}

